\section{A nested-floor sequence and four Pell orbits}\label{sec:pell}

Benoit Cloitre's sequence A081097 lists the positive integers satisfying
\begin{equation}\label{eq:literal-floor}
 n^2=\frac15\left(n+
 \left\lfloor\sqrt5\,n\left\lfloor\sqrt5\,n\right\rfloor\right\rfloor\right).
\end{equation}
The entry records Colin Barker's 2013 empirical generating function
\cite{oeis081097}. We prove that exact formula. The separate conjecture
in the entry concerning arbitrary nonsquare parameters is outside this theorem.
Zero satisfies the displayed equation trivially; the positive restriction
agrees with the source's list and offset.

\begin{theorem}\label{thm:pell-main}
Let $a(1)<a(2)<\cdots$ enumerate the positive solutions of
\eqref{eq:literal-floor}. Put
\[
 U=\begin{pmatrix}161&360\\72&161\end{pmatrix},\qquad
 s_0=\binom42,\quad s_1=\binom{29}{13},\quad
 s_2=\binom{114}{51},\quad s_3=\binom{199}{89}.
\]
For $k\ge0$ and $0\le i<4$,
\[
 a(4k+i+1)=\frac{(U^ks_i)_1+1}{5}.
\]
In particular, the first eight terms are
$1,6,23,40,273,1870,7343,12816$, and
\begin{equation}\label{eq:pell-recurrence}
 a(j+8)=322a(j+4)-a(j)-64\qquad(j\ge1).
\end{equation}
As an identity of formal power series,
\begin{equation}\label{eq:pell-gf}
\sum_{j\ge1}a(j)z^j=
\frac{z(z^7+z^6+13z^5+89z^4-17z^3-17z^2-5z-1)}
{(z-1)(z^2-4z-1)(z^2+4z-1)(z^4+18z^2+1)}.
\end{equation}
\end{theorem}

\begin{lemma}[Exact floor reduction]\label{lem:floor-bridge}
For $n\ge1$, equation \eqref{eq:literal-floor} holds if and only if
there is an integer $y\ge0$ satisfying
\[
 y^2=5n^2-2n+1\quad\hbox{or}\quad y^2=5n^2-2n+2.
\]
Equivalently, with $x=5n-1$,
\begin{equation}\label{eq:pell-admissible}
 x^2-5y^2\in\{-4,-9\},\qquad x\equiv4\pmod5,\qquad x,y\ge0.
\end{equation}
\end{lemma}

\begin{proof}
Suppose first that the nested-floor equation holds, and put
$y=\lfloor\sqrt5\,n\rfloor$, $A=5n^2-n$, and $u=\sqrt5\,ny$.
Then $A\le u<A+1$, with $A,u\ge0$. Dividing the lower inequality by $n$
and squaring gives
\[
 5y^2\ge(5n-1)^2=25n^2-10n+1.
\]
Since $y^2$ is integral, $y^2\ge5n^2-2n+1$.
If $y^2\ge5n^2-2n+3$, the upper inequality, squared, would imply
\[
 25n^4-10n^3+15n^2\le5n^2y^2
 <(5n^2-n+1)^2
 =25n^4-10n^3+11n^2-2n+1,
\]
which contradicts $n\ge1$. This proves the two possibilities.

Conversely, suppose $y^2=5n^2-2n+c$ for $c\in\{1,2\}$ and $y\ge0$.
The inequalities
\[
 (\sqrt5\,n-1)^2<5n^2-2n+1\le y^2\le5n^2
\]
give $y\le\sqrt5\,n<y+1$, hence the inner floor is $y$.
Furthermore,
\[
 u^2-A^2=(5c-1)n^2>0,\qquad
 (A+1)^2-u^2=(11-5c)n^2-2n+1.
\]
The second expression is positive for $c=1$. For $c=2$ it is
$(n-1)^2$, also positive unless $n=1$; that exception would require
$y^2=5$, impossible for an integer. Thus $A\le u<A+1$, proving the
outer floor identity. Substituting $x=5n-1$ gives the stated norms.
\end{proof}

\begin{lemma}[Exhaustive descent]\label{lem:pell-descent}
The pairs satisfying \eqref{eq:pell-admissible} are precisely
$U^ks_i$ for $k\ge0$ and $0\le i<4$.
\end{lemma}

\begin{proof}
The determinant of $U$ is $161^2-5\cdot72^2=1$. Both $U$ and
\[
 U^{-1}=\begin{pmatrix}161&-360\\-72&161\end{pmatrix}
\]
preserve the norm $x^2-5y^2$ and the first coordinate modulo $5$.
Each seed satisfies the required conditions, and the forward map preserves
nonnegativity. This proves one inclusion.

For the converse write $x^2-5y^2=-D$, where $D\in\{4,9\}$.
If $x\ge483$, then
\begin{align*}
 (161x)^2-(360y)^2&=x^2-25920D
       \ge483^2-25920\cdot9=9>0,\\
 5\big((161y)^2-(72x)^2\big)&=x^2+25921D>0,\\
 (360y)^2-(160x)^2&=320x^2+25920D>0.
\end{align*}
All compared quantities are nonnegative, so
$U^{-1}(x,y)$ has positive coordinates and its first coordinate is
strictly less than $x$. This is a descent within the same admissible set.
It terminates at $x<483$. At that point
\[
 5y^2=x^2+D\le482^2+9=232333<5\cdot217^2,
\]
so $y<217$. In this finite rectangle the residue and norm tests give
exactly
\[
\begin{array}{c|rrrr}
 x&4&29&114&199\\
 y&2&13&51&89\\
 x^2-5y^2&-4&-4&-9&-4
\end{array}
\]
and no other pair. This finite claim is explicitly checked by Lean's
kernel in \path{pell_reduced_certificate}; a separate integer script
repeats it. Reversing the descent proves the unbounded classification.
\end{proof}

\begin{proof}[Proof of Theorem~\ref{thm:pell-main}]
The floor bridge and descent give all candidates. Coordinatewise,
\[
 s_0<s_1<s_2<s_3<Us_0=\binom{1364}{610}.
\]
Every entry of $U$ is positive, so these strict inequalities persist
after applying $U^k$. The first coordinates therefore occur in exactly
the interleaved order asserted in the theorem. They are congruent to
$4$ modulo $5$, so $(x+1)/5$ is a positive integer. This proves
the exact enumeration, including uniqueness of every index.

The matrix identity $U^2-322U+I=0$ gives
$x_{k+2}=322x_{k+1}-x_k$ on each orbit. Substituting $x_k=5n_k-1$
gives \eqref{eq:pell-recurrence}. Subtracting consecutive instances yields
\[
 a(j+9)=a(j+8)+322a(j+5)-322a(j+4)-a(j+1)+a(j).
\]
Using the initial terms, coefficient comparison gives
\[
 \sum_{j\ge1}a(j)z^j=
 \frac{z+5z^2+17z^3+17z^4-89z^5-13z^6-z^7-z^8}
 {(1-z)(1-322z^4+z^8)}.
\]
Finally,
\[
 (z^2-4z-1)(z^2+4z-1)(z^4+18z^2+1)=z^8-322z^4+1,
\]
which converts this expression to \eqref{eq:pell-gf}. The denominator
has nonzero constant term, so all manipulations are valid formally.
\end{proof}

\paragraph{Formal endpoint and attribution.}
Six Lean modules prove the full statement from the real-floor definition.
The final theorem \path{A081097.full_resolution} combines strict ordering,
positive terms, unique literal membership, the coefficient identification,
and the exact source numerator--denominator identity in $\Z[[X]]$.
The bounded certificate is only the terminating case of an infinite descent.
Cloitre's definition and Barker's formula retain their attribution;
Pell descent and the matrix generating-function method are classical.
The source and prior-work record is retained in the accompanying package.

